\section{Observation equivalence and minimum order}\label{app:realization}
A realization $(R,T,B)$ is minimal when the columns of $T^jB$, $j\ge0$, span its state space and $RT^jv=0$ for every $j\ge0$ implies $v=0$. A sufficiently large block Hankel matrix $(RT^{i+j}B)_{i,j}$ then has rank equal to the state dimension. Any competing realization factors this same matrix through its own state space, proving the minimum-order assertion without a pole-count shortcut.

\begin{lemma}[Finite certificate and pinned similarity]\label{lem:orbit}
Two realizations of dimensions $n,n'$ have the same all-time kernel if and only if their Markov parameters $RT^jB$ agree for $0\le j<n+n'$. Equivalent minimal realizations have the same dimension and a unique invertible similarity. For the resolved pair of Sec.~\ref{sec:model}, it has the form
\begin{equation}
 S=\diag(I_2,U),\quad U\1=\1,\quad Q'=S^{-1}QS.\label{eq:orbit}
\end{equation}
Conversely, every such similarity yielding an admissible generator preserves the complete kernel. Entries of $U$ need not be nonnegative.
\end{lemma}
\begin{proof}
Apply Cayley--Hamilton to $\diag(T,T')$: zero differences of the first $n+n'$ parameters force all subsequent differences to vanish. The convergent exponential series then gives all-time equality; the converse follows by differentiation. For minimal realizations define, on reachable vectors,
\begin{equation}
 S\left(\sum_jT'^jB'v_j\right)=\sum_jT^jBv_j.
\end{equation}
Equality of parameters and observability make this well defined; reversing the argument gives its inverse. Controllability gives $SB'=B$, $ST'=TS$, and $RS=R'$, and uniqueness. The zero-time data give $q_{xy}=\psi_{-,+}(0)$ and $q_{yx}=\psi_{+,-}(0)$. The common $R,B$ therefore pin every visible row and column of $S$, giving its block form. Since $T\1=-B\1$ and $T$ is invertible, $TS\1=ST'\1=-SB\1=T\1$ implies $S\1=\1$. The observed matrices act only on the fixed visible coordinates, hence $Q'=S^{-1}QS$. Direct substitution proves the converse.
\end{proof}
This is the realization-basis argument specialized to the event maps; its positive-model antecedents are discussed in Sec.~I~\cite{Vanluyten2008,HorvathTelek2007}.

For the balanced source, the hidden blocks are
\begin{equation}
 X=\begin{pmatrix}2&0&8\\0&7&11\end{pmatrix},\quad
 Y=\begin{pmatrix}3&0\\0&6\\z&z\end{pmatrix},\quad
 H=\begin{pmatrix}-7&4&0\\5&-11&0\\0&0&-2z\end{pmatrix}.\label{eq:blocks}
\end{equation}
Writing $Y_j$ for its columns, direct determinants give
\begin{align}
 \det(Y_1,Y_2,HY_1)&=-9z(4z-9),\\
 \det(Y_1,Y_2,HY_2)&=-18z(2z-3),\\
 \det\begin{pmatrix}X_{1,:}\\X_{2,:}\\(XH)_{2,:}\end{pmatrix}&=-14(22z+19).
\end{align}
For $z>0$ the first two cannot both vanish and the last never vanishes. Thus the hidden transfer has reachable and observable order three. The visible plane is directly selected by $R$ and $B$. Subtracting its contributions in successive powers of $T$ leaves the hidden reachability and observability spaces just computed, so the full order is five. This includes both hidden-pole coincidences and $z=2$. The cap forces every competitor to have exactly five states and be minimal.

The observations also determine the visible block $A$: their zero-time values give the observed rates and first derivatives give the killed visible diagonals. Write $A_{\mathrm{kill}}=A-(E_++E_-)_{VV}$. After row permutation and division by observed rates they determine the visible resolvent $K(\lambda)$. Its Schur complement determines
\begin{equation}
 X(\lambda I-H)^{-1}Y=\lambda I-A_{\mathrm{kill}}-K(\lambda)^{-1}.\label{eq:htransfer}
\end{equation}
This is a matrix-valued constraint shared by all competitors. Treating its entries as unrelated scalar phase-type fits would lose it.

\section{Global support exhaustion and the invariant triangle}\label{app:cone}
\subsection{Modal data and disconnected hidden supports}
Let $\rho=\sqrt6$, $\alpha=9-2\rho$, $\beta=9+2\rho$, $r_y=4(\rho-1)/5$, and set
\begin{gather}
 D=\diag(-\alpha,-\beta,-2z),\quad X_m=(p,q,v),\quad
 Y_m=\begin{pmatrix}r\\f\\w\end{pmatrix},\label{eq:modal}\\
 p=\binom{3+\rho/2}{35\rho/8},\quad
 q=\binom{3-\rho/2}{-35\rho/8},\quad v=\binom8{11},\\
 r=(1,r_y),\quad f=(1,-4(\rho+1)/5),\quad w=(z,z).
\end{gather}
Direct multiplication yields
\begin{equation}
 X(\lambda I-H)^{-1}Y
 =\frac{pr}{\lambda+\alpha}+\frac{qf}{\lambda+\beta}+\frac{vw}{\lambda+2z}.\label{eq:residues}
\end{equation}
The residues $pr,vw$ are strictly positive and rank one; $qf$ has negative off-diagonal entries. At $z\ge\zs$ the three poles are distinct and $2z>\beta$.

A disconnected bidirected hidden support on three states consists of three singletons or a pair and a singleton. Singletons have nonnegative outer-product residues, excluding three singletons and a singleton at $\beta$. A singleton at $\alpha$ leaves a pair with dominant eigenvalue $-\beta$; its Perron transfer residue is nonnegative, again impossible. The singleton must carry $-2z$.

The remaining pair has zeroth response $pr+qf=\diag(6,42)$. Nonnegativity and reciprocal visible incidence force its two states to attach separately to $x$ and $y$: a state attached to both would give positive off-diagonal response. Write entrances $x_1,x_2$, exits $y_1,y_2$, and internal rates $h_0,k_0$. Zeroth and first moments, with row sums, give
\begin{gather}
 x_1y_1=6,\quad x_2y_2=42,\quad y_1+h_0=7,\quad y_2+k_0=11,\\
 x_1h_0y_2=48,\qquad x_2k_0y_1=105.
\end{gather}
Hence $h_0k_0=20$ and $19h_0-h_0k_0=56$, giving $h_0=4$, $k_0=5$, $y_1=3$, $y_2=6$, $x_1=2$, $x_2=7$. The singleton's residue and exit sum give entrances $(8,11)^T$ and exits $(z,z)$. Thus this branch contains only the source.

\subsection{All connected hidden supports give finite triangles}
Throughout the connected-case argument take $z\ge21/2$, so $2z>\beta$ and the modal spectrum is distinct. For a connected hidden support, minimality gives an invertible real $V$ with $H'=V^{-1}DV$ Metzler irreducible, $X'=X_mV\ge0$, and $Y'=V^{-1}Y_m\ge0$. The right Perron eigenvector $V^{-1}e_1$ at $-\alpha$ is positive: its sign is fixed by $X'V^{-1}e_1=p>0$. The left eigenvector $e_1^TV$ is then positive as well, since its product with that right eigenvector is one. Every column of $V$ has positive first coordinate. Sectioning $V\mathbb R_+^3$ at that coordinate equal to one gives a finite nondegenerate triangle with the origin strictly inside. Positively rescale its columns to $(1,u_i,v_i)^T$.

After rescaling time by $\beta-\alpha$, the section is invariant under
\begin{equation}
 (u,v)\mapsto(e^{-t}u,e^{-kt}v),\qquad
 k=\frac{2z-9+2\rho}{4\rho}.\label{eq:flow}
\end{equation}
It contains inputs $(1,A)$ and $(-m,B_0)$ and lies in an output wedge, where
\begin{gather}
 A=z,\quad B_0=z/r_y,\quad m=(7+2\rho)/5,\\
 L=(3-\rho/2)/8,\quad R_w=35\rho/88,\quad
 v\ge L(-m-u),\quad v\ge R_w(u-1).\label{eq:wedge}
\end{gather}
These are direct consequences of $Y'\ge0$, $X'\ge0$, and semigroup positivity. No sign constraint on $V$ itself was imposed. A ray with zero first coordinate is impossible in this irreducible branch; the disconnected branch was already treated separately.

Any wedge point of nonpositive height has abscissa between $-m$ and $1$. A triangle with only one positive-height vertex could not contain both positive-height inputs: the same vertex would have to extend beyond both ends of that interval. The interior origin forces a negative-height vertex. Thus the vertices can be labeled
\begin{equation}
 (-\ell,P),\quad(c,Q_v),\quad(d,-W),\qquad
 \ell\ge m,\ c\ge1,\ -m<d<1,\ P,Q_v,W>0.\label{eq:vertices}
\end{equation}
Put $a=k-1$. The output, lower-face invariance, upper containment, and lower containment conditions include
\begin{align}
 &P\ge L(\ell-m),\quad W\le\min\{L(m+d),R_w(1-d)\}<1,\label{eq:ineq1}\\
 &(a\ell+kd)P\le\ell W,\qquad(ac-kd)Q_v\le cW,\label{eq:ineq2}\\
 &B_0\le H_c:=\frac{(c+m)P+(\ell-m)Q_v}{\ell+c},\label{eq:height}\\
 &A\le\frac{(c-1)P+(\ell+1)Q_v}{\ell+c},\label{eq:upperA}\\
 &(Q_v+W)(1-d)\le(A+W)(c-d),\label{eq:lowerA}\\
 &(P+W)(d+m)\le(B_0+W)(\ell+d).\label{eq:lowerB}
\end{align}
For example, differentiating the lower-left supporting line along~\eqref{eq:flow} at $(-\ell,P)$ gives the first inequality~\eqref{eq:ineq2}. All displayed geometric denominators are positive.

We now restrict to $z\ge21/2$, sufficient for the threshold and above. If $d<0$, the reflected use of~\eqref{eq:ineq1},~\eqref{eq:ineq2},~\eqref{eq:upperA}, and~\eqref{eq:lowerB} gives
\begin{equation}
 A<\frac{L}{a}\left[m+\frac{B_0+1}{R_w}\right]
 <\frac{25/8+(25/18)(z+1)}{z-7}.
\end{equation}
The latter is smaller than $z$ for $z\ge10$, because its reversed difference has numerator
$72(z-10)^2+836(z-10)+835>0$. Here we used $m<5/2$, $L<1/4$, $R_w>9/10$, $B_0<z$, and $a>(z-7)/5$. Thus $d\ge0$. Now $a>7/10$ and~\eqref{eq:ineq2} give $P<10/7<2<B_0<A$. Equations~\eqref{eq:height}--\eqref{eq:upperA} imply $Q_v\ge A$ and $\ell>m$.

Write $\epsilon=1-d$. Since $aQ_v-W>0$, the second inequality~\eqref{eq:ineq2} and $c\ge1$ imply
\begin{equation}
 \epsilon\le\frac{A+W}{kA}\le\frac{A+R_w\epsilon}{kA},\qquad
 \epsilon\le\frac1{k-R_w/A}.\label{eq:apexbound}
\end{equation}
The wedge apex has
\begin{equation}
 d_0=\frac{R_w-Lm}{L+R_w}>0,\qquad
 1-d_0=\frac{L(m+1)}{L+R_w}.
\end{equation}
At $z=21/2$, the last fraction minus the bound in~\eqref{eq:apexbound} exceeds $0.014259$ by rational evaluation with $\rho^2=6$ and $\rho>0$. Since $k-R_w/z$ increases, $d>d_0$ throughout the range. Therefore the active bottom constraint is $W\le R_w\epsilon$.

\subsection{Global envelope and its unique equality case}
At fixed $d,W$, the left inequalities imply
\begin{equation}
 L(\ell-m)\le g(\ell):=\frac{\ell W}{a\ell+kd}.
\end{equation}
Since $d>0$, $g$ increases strictly and is concave. There is one intersection $\ell_*=m+t$, $t>0$, with
\begin{equation}
 P_*=Lt=g(m+t),\qquad Lt[a(m+t)+kd]=(m+t)W.\label{eq:leftmax}
\end{equation}
The corresponding quadratic in $t$ has positive leading coefficient and negative constant. All admissible $\ell,P$ satisfy $\ell\le\ell_*$, $P\le P_*$. In~\eqref{eq:height}, both partial derivatives in $P$ and $\ell$ are strictly positive when $Q_v>P$; explicitly they are $(c+m)/(\ell+c)$ and $(c+m)(Q_v-P)/(\ell+c)^2$. Replacing the left vertex by its maximum therefore bounds every feasible triangle, even if a relaxed triangle violates an unused constraint.

For the right vertex,~\eqref{eq:lowerA} bounds $Q_v$ by an increasing line,
$Q_v\le (A+W)(c-d)/\epsilon-W$. On $ac-kd>0$,~\eqref{eq:ineq2} gives the decreasing hyperbola $Q_v\le cW/(ac-kd)$. The relevant intersection is
\begin{equation}
 c_* =\frac{k}{a}\left[d+\frac{W\epsilon}{A+W}\right],\qquad
 Q_*^v=\frac{Ad+W}{a\epsilon}.\label{eq:rightmax}
\end{equation}
It lies strictly beyond the pole $kd/a$. The other algebraic intersection is $c=d<1$ and is outside the domain. If $c_*<1$, the hyperbola at $c=1$ is below $A$, contradicting $Q_v\ge A$. Along the line, $H_c$ increases in $c$, with derivative numerator $(\ell_*-m)[(A+W)(\ell_*+d)/\epsilon+W+P_*]>0$. Along the hyperbola it decreases whenever $Q_v>P_*$; a point with $Q_v\le P_*$ cannot attain $B_0>2$. Thus the unique maximal right choice is~\eqref{eq:rightmax}.

Set $r_c=c_*/Q_*^v=k\epsilon/(A+W)$. The resulting envelope is
\begin{equation}
 H_* =\frac{t[Lm+(1+Lr_c)Q_*^v]}{m+t+r_cQ_*^v}.\label{eq:envelope}
\end{equation}
As $W$ increases, $t$ and $Q_*^v$ increase strictly and $r_c$ decreases strictly. The partial derivatives of~\eqref{eq:envelope} with respect to $t,Q_*^v$ are positive; its derivative in $r_c$ has sign $Lt-Q_*^v<0$. These signs persist up to $W=R_w\epsilon$, since $Lt<2<A\le Q_*^v$. Hence the final maximum uses $W=R_w\epsilon$.

Let $J=am+k$. Equations~\eqref{eq:leftmax} and $W=R_w\epsilon$ give
\begin{equation}
 \epsilon=\frac{Lt(J+at)}{R_wm+(kL+R_w)t},\qquad
 H_*-B_0=\frac{f_0+f_1t+f_2t^2}{(J+at)^2},\label{eq:quadratic}
\end{equation}
where
\begin{align}
 f_0&=AR_wm/L-B_0J^2,\\
 f_1&=A(R_w/L+k-am)+LJ^2+R_wJ-2B_0Ja,\\
 f_2&=a[LJ+R_w-A(1+Lk/R_w)-B_0a].\label{eq:fcoeff}
\end{align}
These identities follow by direct rational substitution, with no optimization assumptions. The coefficients obey
\begin{align}
 f_0&=-\frac{19+9\rho}{100}z\left(z^2-11z+\frac{211}{44}\right),\label{eq:f0}\\
 f_1^2-4f_0f_2
 &=-\left(\frac{11}{1200}+\frac{11\rho}{4200}\right)
 \left(z+\frac{19}{22}\right)^2\frac{\mathcal P(z)}{352}.\label{eq:discriminant}
\end{align}
Descartes' rule, $\mathcal P(0)<0$, and $\mathcal P(z)\to+\infty$ as $z\to+\infty$ give a unique positive root. Rational evaluation isolates it in $(10.55714966866,10.55714966867)$. The larger root of the quadratic in~\eqref{eq:f0} is below this interval. Thus for $z>\zs$, $f_0<0$ and the discriminant is negative; the quadratic in~\eqref{eq:quadratic} is strictly negative everywhere. This excludes every connected support.

At $z=\zs$, exact root-interval evaluation gives
\begin{equation}
 -0.520929<f_0<-0.520928,\quad
 5.496508<f_1<5.496509,\quad
 -14.498918<f_2<-14.498917.
\end{equation}
The quadratic is nonpositive and has its sole zero at $t_*=-f_1/(2f_2)>0$. Feasibility requires equality throughout the chain $B_0\le H_c\le H_*\le B_0$. Strictness of every maximization fixes the unique vertices through~\eqref{eq:leftmax},~\eqref{eq:rightmax}, and~\eqref{eq:quadratic}. The left and bottom vertices lie on the left and right output faces, respectively. Input $(1,A)$ lies on the lower-right edge, input $(-m,B_0)$ on the upper edge, and both upper-vertex lower-face derivatives vanish. These paired input/output zeros eliminate one visible pair at each of two vertices; the two tangencies eliminate both directions between the upper vertices. The resulting support is precisely the hidden path specified below.

The unique triangle gives at most one normalized generator, not merely one cone. For its vertex matrix $V$, put $\bar H=V^{-1}DV$, $\bar Y=V^{-1}Y_m$, $\bar X=X_mV$. Column rescaling is fixed uniquely by
\begin{equation}
 w=-\bar H^{-1}\bar Y\1,\quad W_d=\diag(w),\quad
 H'=W_d^{-1}\bar HW_d,\quad X'=\bar XW_d,\quad Y'=W_d^{-1}\bar Y.\label{eq:normalize}
\end{equation}
All hidden eigenvalues are nonzero. Existence, positive scaling, and entropy separation are supplied independently in Appendix~\ref{app:certificate}; they do not follow merely from equality in a necessary bound.

\section{Exact endpoint construction and entropy certificate}\label{app:certificate}
An explicit rational definition avoids using rounded rates. In the following equations evaluate all expressions at $z=\zs$. Define
\begin{equation}
 C=\begin{pmatrix}3&0&1\\0&6&1\\z&z&1\end{pmatrix},\quad D_c=C^{-1}HC,
\end{equation}
where $H$ is~\eqref{eq:blocks}. Use one-based indices and set
\begin{equation}
 (a_1,a_2,b_1,b_2,c_1,c_2)
 =((D_c)_{22},(D_c)_{21},(D_c)_{12},(D_c)_{11},(D_c)_{32},(D_c)_{31}).
\end{equation}
Let
\begin{align}
 a_p(z)&=100z^4-840z^3+3741z^2-6867z+3978,\\
 b_p(z)&=-408z^4-2916z^3+8766z^2+3105z-12474,\\
 c_p(z)&=2412z^4-5724z^3+891z^2+3402z.
\end{align}
They satisfy the independent physical-path discriminant identity
\begin{equation}
 b_p^2-4a_pc_p=-567(z-2)^2(2z-3)^2\mathcal P(z).\label{eq:pathdisc}
\end{equation}
Set
\begin{gather}
 u=-\frac{b_p}{2a_p},\qquad v=-\frac{u(a_2u+c_2)}{u^2+a_1u+c_1},\\
 p_x=(42+11z)/18,\quad p_y=11z/18,\quad
 q_x=8z/10,\quad q_y=(6+8z)/10,\\
 c_p'=\frac{p_x-u}{p_y},\qquad d_p'=\frac{q_y-v}{q_x},\qquad
 A_p=\frac{u+c_p'v}{1-c_p'd_p'},\quad
 B_p=\frac{v+d_p'u}{1-c_p'd_p'}.\label{eq:pathcoordinates}
\end{gather}
The primed $c_p',d_p'$ are chart coordinates, not the polynomial $c_p(z)$. Finally define
\begin{equation}
 C'=\begin{pmatrix}0&u&1\\v&0&1\\B_p&A_p&1\end{pmatrix},\quad
 U=C(C')^{-1},\quad S=\diag(I_2,U),\quad Q_*=S^{-1}Q(\zs)S.\label{eq:target}
\end{equation}
The hidden order is $(a,b,c)$: a $y$-only endpoint, an $x$-only center, and a shared endpoint. The seven reciprocal edges are $xy,xb,xc,ya,yc,ab,bc$. All other off-diagonals are exactly zero. For orientation,
\begin{equation}
 Q_*\simeq\begin{pmatrix}
 -11&1&0&1.86553087&8.13446913\\
 2&-20&6.92248992&0&11.07751008\\
 0&6.22824145&-11.27299011&5.04474866&0\\
 2.77737298&0&3.73491843&-6.73129745&0.21900604\\
 10.48328058&10.38263173&0&0.24409946&-21.11001177
 \end{pmatrix}.
\end{equation}
This rounded matrix is not the equality certificate.

The following is a finite, reproducible algebraic admissibility check. Work in the real field $\mathbb Q(\zs)$ with $\mathcal P(\zs)=0$, reducing each element to a polynomial of degree at most three. Isolate the positive root in
\begin{equation}
 \left[\frac{6329583361674357233}{599554194108037085},
 \frac{5875360061222171594}{556529010729195633}\right].\label{eq:rootinterval}
\end{equation}
The interval has width below $10^{-35}$; opposite polynomial signs and unique positivity select the embedding. Rational interval evaluation verifies that $\det C=18-9z$, $a_p$, $u^2+a_1u+c_1$, $p_y,q_x$, $1-c_p'd_p'$, and $\det C'$ are all nonzero. Reduce~\eqref{eq:target} exactly: $U\1=\1$, row sums vanish, the six absent off-diagonal entries are zero, and the fourteen remaining off-diagonals are positive. Their enclosures lie strictly between $1/5$ and $12$. The stated seven-edge graph is connected. This suffices for admissibility. The package supplies every reduced coefficient, making these finite checks repeatable without solving a new realization problem.

Compute the unique normalized stationary vectors by exact reduced linear solves. Both are positive and satisfy $\pi_* =\pi S$. Directly verify $RT^jB=RT_*^jB$ for $j=0,\ldots,9$; Lemma~\ref{lem:orbit} gives all-time equality independently of the displayed similarity. As a separate coordinate check, the equality triangle in Appendix~\ref{app:cone}, normalized by~\eqref{eq:normalize}, agrees coefficientwise with~\eqref{eq:target} after swapping its last two hidden vertices. This computation uses pairs $a+b\sqrt6$ over $\mathbb Q(\zs)$; every final $\sqrt6$ coefficient cancels.

Entropy separation also needs a certificate, since different generators need not have different entropy. Evaluate stationary currents and positive flux ratios with rational intervals from~\eqref{eq:rootinterval}. To bound a logarithm, reduce its positive rational endpoints by powers of two into $[1,2]$. For $t=(x-1)/(x+1)\in[0,1/3]$ use
\begin{equation}
 \log x=2\sum_{j=0}^{n-1}\frac{t^{2j+1}}{2j+1}+R_n,\qquad
 0\le R_n\le\frac{2t^{2n+1}}{(2n+1)(1-t^2)},\quad n=24.\label{eq:logbound}
\end{equation}
Bound $\log2$ by the same formula. Interval arithmetic accounts for the signs of powers-of-two shifts, currents, and products, and sums once per unordered present edge. It gives~\eqref{eq:sig1}--\eqref{eq:sig2} and
\begin{equation}
 0.002633285729\le\sig(Q_*)-\sig(Q(\zs))\le0.002633285730.
\end{equation}
These are rational enclosures, not precision-based equality tests. Together with Appendix~\ref{app:cone}, they finish the endpoint and supercritical parts of Theorem~\ref{thm:balanced}.

\section{An exact entropy-divergent hidden-pair orbit}\label{app:pair}
Take fully hidden states $h,k$ with the same external neighbor set $J$, and different full external exit rows $Y_h,Y_k$. Let $m_0=q_{hk}$, $n_0=q_{kh}$, and let
\begin{equation}
 \lambda_h=m_0+\sum_{i\in J}Y_{hi},\qquad
 \lambda_k=n_0+\sum_{i\in J}Y_{ki}
\end{equation}
be their \emph{total} escape rates. Assume either $m_0,n_0>0$, or $m_0=n_0=0$ with $\lambda_h=\lambda_k$. The external rows include all other hidden states, not only visible endpoints.

Exchange $h,k$ if necessary so that $r_*:=\min_{i\in J}Y_{hi}/Y_{ki}\in(0,1)$. Such an orientation exists because the rows differ and their supported entries are positive. Embed
\begin{equation}
 U_e=\begin{pmatrix}1-e&e\\0&1\end{pmatrix},\qquad 0\le e<r_*,
\end{equation}
as the hidden block of $S_e$. The transformed outside entrances and exits are
\begin{align}
 q'_{ih}&=(1-e)q_{ih},&q'_{ik}&=q_{ik}+eq_{ih},\\
 q'_{hi}&=(Y_{hi}-eY_{ki})/(1-e),&q'_{ki}&=Y_{ki}.
\end{align}
They retain their reciprocal support before $r_*$. For an adjacent pair,
\begin{equation}
 q'_{kh}=n_0(1-e),\qquad
 q'_{hk}=\frac{f(e)}{1-e},\quad
 f(e)=m_0+e(\lambda_k-\lambda_h)-n_0e^2.\label{eq:shear}
\end{equation}
Let $e_*$ be the first zero of $f$ in $(0,r_*]$, if present, and otherwise $r_*$. Then $0<e_*\le r_*<1$, and all generators before $e_*$ retain the original support. For the unjoined equal-escape pair the hidden block is scalar and unchanged; take $e_*=r_*$.

Normalized similarity preserves the exact kernel. The stationary law is explicitly
\begin{equation}
 \pi'_i=\pi_i\ (i\notin\{h,k\}),\quad
 \pi'_h=(1-e)\pi_h,\quad\pi'_k=\pi_k+e\pi_h.\label{eq:pipair}
\end{equation}
Every stationary mass has a positive finite limit and all rates remain bounded because $e_*<1$. If $f(e_*)=0$, the $h\to k$ flux vanishes while its reverse tends to $(\pi_k+e_*\pi_h)n_0(1-e_*)>0$. If $e_*=r_*$, an external exit $h\to i$ vanishes while its reverse tends to $\pi_i(1-r_*)q_{ih}>0$. This also covers simultaneous zeros. The corresponding nonnegative entropy term diverges, and other terms cannot cancel it. The endpoint itself need not be admitted or have a unique stationary law;~\eqref{eq:pipair} gives the necessary limits before the boundary.

For a minimal realization, equal full external exit rows are impossible: the codimension-one subspace of vectors with equal $h,k$ entries is invariant under $T$ and contains every column of $B$. Reachability would fail. Thus every pair in a complete minimal realization meets the unequal-row requirement. This observation is used below, while Lemma~\ref{lem:neighborhood} creates unequal rows explicitly and does not assume minimality.

\section{Complete subcritical coverage, including repeated poles}\label{app:subcritical}
We give four covering constructions. Each preserves the hidden transfer exactly and yields a pair to which Appendix~\ref{app:pair} applies.

\subsection{Below the slow pole}
Take $0<z<\alpha/2$. Let $d_s=2(\rho-1)$. In hidden order $(h,k,l)$ set
\begin{equation}
 K_{hl}=-4,\quad K_{kl}=-d_s,\quad K_{lh}=5,\quad K_{lk}=d_s,\quad
 K_{hk}=K_{kh}=0,
\end{equation}
and fix diagonals by $K\1=0$. Under $U=I+\epsilon K$, the derivative of the transformed generator is the commutator with embedded $K$. The four missing hidden rates have derivatives $(\alpha-2z)(4,d_s,5,d_s)>0$; the missing visible incidences $(xk,yh,hy,kx)$ have derivatives $(8d_s,55,4z,d_sz)>0$. For sufficiently small $\epsilon>0$, $U$ is invertible, all missing pairs open, and existing positive rates remain positive. The target is complete and minimal, hence has an entropy-divergent hidden pair. This interval includes the path-chart singularity $z=2$.

\subsection{Between the hidden poles}
Take $\alpha/2\le z\le\beta/2$. Put $s_0=2z$, $d_s=2(\rho-1)$, and $\Delta=4\rho$. Factor the two relevant residues as
\begin{equation}
 Z_{s_0}=\binom8{11}(z,z),\qquad
 Z_\beta=\binom{2d_s}{-35}\frac{(3,-24/d_s)}{\Delta}.
\end{equation}
Choose midpoints in the nonempty intervals
\begin{equation}
 \frac{24}{d_s\Delta z}<t<\frac4{d_s},\qquad
 \frac3{\Delta z}<u<\frac{11}{35}.\label{eq:twomodeintervals}
\end{equation}
They are nonempty because $z>\rho/4$ and $z>35\rho/88$, both already true at $\alpha/2$. For $P_2=\left(\begin{smallmatrix}1&1\\-t&u\end{smallmatrix}\right)$, use
\begin{align}
 \bar X&=\begin{pmatrix}8-2d_st&8+2d_su\\11+35t&11-35u\end{pmatrix}>0,\\
 \bar Y&=\frac1{t+u}\begin{pmatrix}uz-3/\Delta&uz+24/(d_s\Delta)\\tz+3/\Delta&tz-24/(d_s\Delta)\end{pmatrix}>0,\\
 \bar H&=\frac1{t+u}\begin{pmatrix}-s_0u-\beta t&u(\beta-s_0)\\t(\beta-s_0)&-s_0t-\beta u\end{pmatrix}.
\end{align}
These are obtained by similarity from modes $-s_0,-\beta$. The stable Metzler $\bar H$ has two positive off-diagonals for $s_0<\beta$, and is $-\beta I$ at equality. Normalize by~\eqref{eq:normalize}, giving a pair with strictly positive visible incidences. Its exit matrix has rank two because the modal exit rows are a positive row and a mixed-sign row; their independence is preserved.

For the remaining residue write
\begin{equation}
 Z_\alpha=cg,\quad c=\binom{2(d_s+4)}{35},\quad
 g=(3,24/(d_s+4))/\Delta.
\end{equation}
A singleton with entrance $(g\1)c/\alpha$ and exit $\alpha g/(g\1)$ realizes it. Adjoin it without hidden links to the pair, retaining the source visible block. Hidden row normalization and equality of the transfer at zero give the source's visible hidden-entry totals $(10,18)^T$, so visible rows also normalize. The full generator is irreducible through its visible links. The pair has unequal external exits, is adjacent for $s_0<\beta$, and unjoined with equal escapes at $s_0=\beta$. Appendix~\ref{app:pair} applies in both cases. All transfer identities remain valid when poles coincide: the residues simply add. Thus both endpoints, including $s_0=\alpha$, are covered without a distinct-pole assumption.

\subsection{An explicit fast opening}
Take $\beta/2\le z\le21/2$. Put $M_s=2z-11>0$, $L_s=M_s+4$, and define
\begin{gather}
 v_{max}=7z/33,\quad
 q=\tfrac12\left(M_s/5+96v_{max}/(zM_s)\right),\quad
 v=\tfrac12\left(zqM_s/96+v_{max}\right),\\
 u=\tfrac12(zq/24+4v/M_s),\quad
 p=\tfrac12(-4u-zq/6),\quad
 r=\tfrac12(11v/7+z/3).
\end{gather}
Use hidden off-diagonals $(K_{hk},K_{hl},K_{kh},K_{kl},K_{lh},K_{lk})=(p,q,r,-1,-v,u)$ and $K\1=0$. Since $L_sM_s\ge20$ and $11M_s^2\le1100<1120$, the midpoint intervals are strict:
\begin{gather}
 q>M_s/5,\quad v<7z/33,\quad zq/24<u<4v/M_s,\\
 -4u<p<-zq/6,\qquad11v/7<r<z/3.
\end{gather}
The eight missing-rate derivatives, ordered $(xk,yh,hy,hl,kx,kl,lh,lk)$, are
\begin{equation}
 (2p+8u,\ 7r-11v,\ -6p-zq,\ L_sq-4,\ -3r+z,\ 5q-M_s,\ L_sv-5u,\ 4v-M_su),
\end{equation}
all positive. Thus $I+\epsilon K$ opens a complete realization for small $\epsilon>0$, with exact observation and subsequent entropy divergence as above. This argument needs neither an exit chart nor an assumption that failure of a first-order opening implies global rigidity.

\subsection{Strict triangles up to the fold}
Take $21/2\le z<\zs$. Use the quantities of Appendix~\ref{app:cone} and set $t=-f_1/(2f_2)$. On the entire rational interval
\begin{equation}
 I=[10.5,10.55714966867],\label{eq:wholeinterval}
\end{equation}
rational interval evaluation gives $f_1>5$ and $f_2<-14$. For $z<\zs$, the positive discriminant in~\eqref{eq:discriminant} gives
\begin{equation}
 f_0+f_1t+f_2t^2=-\frac{f_1^2-4f_0f_2}{4f_2}>0.
\end{equation}
Define $\epsilon,d,W,\ell,P,c,Q_v$ by~\eqref{eq:quadratic}, $d=1-\epsilon$, $W=R_w\epsilon$, $\ell=m+t$, $P=Lt$, and~\eqref{eq:rightmax}. Then $B_0<H_*$ is strict upper-input slack. The remaining geometric signs hold throughout $I$; Table~\ref{tab:margins} gives conservative rational lower bounds. To reproduce them, enclose $\rho$ between the exact decimal rationals $2.449489742783178098197284$ and $2.449489742783178098197285$, whose squares straddle six. Shift polynomial coefficients by $z=21/2+\xi$, and apply rational interval operations, rounded outward to 24 decimal places after each operation, with $0\le\xi\le0.05714966867$. No parameter sampling establishes these bounds.

\begin{table}[ht]
\caption{Strict rational lower bounds valid on the whole interval~\eqref{eq:wholeinterval}, with $p_\ell=(P+W)/(\ell+d)$, $q_r=(Q_v+W)/(c-d)$. The upper-input-$B_0$ slack is not listed because it vanishes at the fold.}\label{tab:margins}
\begin{ruledtabular}
\begin{tabular}{lc@{\qquad}lc}
Quantity & Lower bound & Quantity & Lower bound\\
\hline
$t$ & $18/100$ & $\epsilon$ & $5/100$\\
$1-\epsilon$ & $9/10$ & $c-1$ & $1$\\
$L(m+d)-W$ & $6/10$ & $Q_v-R_w(c-1)$ & $170$\\
$B_0+W-p_\ell(d+m)$ & $8$ & $[(c-1)P+(\ell+1)Q_v]/(\ell+c)-A$ & $100$\\
$W-p_\ell d$ & $3/100$ & $W+q_r d$ & $120$\\
$kW-p_\ell d$ & $7/100$ & $kW+q_r d$ & $120$
\end{tabular}
\end{ruledtabular}
\end{table}

For completeness, the left and right lower-face tangencies imply $p_\ell\ell=kP$ and $q_rc=kQ_v$. Therefore the lower-face margins of the origin are $W-p_\ell d=aP>0$ and $W+q_rd=aQ_v>0$. The upper edge is above the origin, so the triangle is nondegenerate with an interior origin. The inward derivatives at the bottom vertex are $a(W+P)>0$ and $a(W+Q_v)>0$. The upper-edge inward derivatives at its endpoints have positive numerators $[a\ell+kc]P+\ell Q_v$ and $[k\ell+ac]Q_v+cP$. The two remaining derivatives vanish by tangency. These vertex-face conditions suffice for invariance under the linear flow.

The initial triangle is nonnegative but can have one-way support. Make three small changes for each fixed $z<\zs$. First decrease $W$ while holding $d$ fixed, recomputing~\eqref{eq:leftmax} and~\eqref{eq:rightmax}; all strict margins persist and the bottom output becomes strict. Next decrease $\ell$ slightly and take $P$ to be the midpoint of $L(\ell-m)$ and $\ell W/(a\ell+kd)$; both left output and left inward inequalities become strict. Finally keep $c=c_*$ and decrease $Q_v$ slightly; since $ac-kd>0$, both the right inward inequality and the lower containment of $(1,A)$ become strict. All other strict inequalities persist by continuity. Only finitely many conditions are involved, so successive sufficiently small positive choices exist; no uniform size as $z\uparrow\zs$ is claimed.

The final triangle has strictly positive input barycentric coordinates, output values, and all off-diagonal inward derivatives. At vertex $j$, the derivative of each initially zero barycentric coordinate $i\ne j$ is $(\bar H)_{ij}$ up to the positive time-rescaling factor; thus strict inward derivatives give positive off-diagonals. Consequently $\bar X>0$, $\bar Y>0$, and every off-diagonal of the stable $\bar H$ is positive. The normalization vector in~\eqref{eq:normalize} is positive because $-\bar H^{-1}>0$. Hidden row sums vanish and transfer equality at zero fixes the visible row sums. This is a complete five-state physical generator with the same all-time observation. It is minimal by Appendix~\ref{app:realization}, so its hidden pair meets Appendix~\ref{app:pair}. The four intervals cover all $0<z<\zs$, completing Theorem~\ref{thm:balanced}.

\section{Why one additional state changes the exact answer}\label{app:split}
Let $h$ be fully hidden in any irreducible generator. Replace it by $h_1,h_2$; keep rates among outside states, split each incoming rate as $\theta_iq_{ih}$ and $(1-\theta_i)q_{ih}$ with $0<\theta_i<1$, and give both clones the original outgoing rates $q_{hj}$ to every outside $j$. Set internal rates $q_{h_1h_2}=m>0$ and $q_{h_2h_1}=n>0$, with $n$ fixed. Normalize all rows. Every positive $m$ gives connected reciprocal support with the same observed edges.

For the aggregation matrix $C$ merging the clones, $\widetilde Q C=CQ$. Because no mark touches $h$, also $\widetilde T C=CT$, $\widetilde B=CB$, and $\widetilde R C=R$. Therefore $\widetilde\Psi(t)=\widetilde R e^{\widetilde Tt}CB=Re^{Tt}B$ exactly. This is strong lumping applied to the killed generator, not just agreement of stationary probabilities.

Put $\lambda=\sum_{j\ne h}q_{hj}$ and $\alpha_c=\sum_{i\ne h}\pi_i\theta_iq_{ih}/\pi_h\in(0,\lambda)$. Outside stationary probabilities remain $\pi_i$, and clone probabilities are
\begin{equation}
 \widetilde\pi_{h_1}=\pi_h\frac{\alpha_c+n}{\lambda+m+n},\qquad
 \widetilde\pi_{h_2}=\pi_h\frac{\lambda+m-\alpha_c}{\lambda+m+n}.
\end{equation}
Both have positive limits as $m\downarrow0$. The first internal flux tends to zero and the reverse tends to $\pi_hn(\lambda-\alpha_c)/(\lambda+n)>0$, so entropy diverges while rates remain uniformly bounded above. Applied to $Q(\zs)$, this proves unboundedness with six states. These enlarged realizations are nonminimal for the common marked kernel. This appendix asserts sufficiency of one extra state, not a general necessary state budget.


\section{Local persistence of the five-state boundary}\label{app:unfolding}
\begin{proof}
In Appendix~\ref{app:cone} replace only $A$ and $B_0$ by
\begin{equation}
 A=z+\delta,\qquad B_0=(z-\delta)/r_y.
\end{equation}
The other modal constants, the flow, and the wedge are unchanged.
Minimality persists: one hidden reachability minor and the observability
minor are
\begin{align}
 \det(Y_1,Y_2,HY_1)&=-9(4\delta z-19\delta+4z^2-9z),\\
 \det\begin{pmatrix}X_{1,:}\\X_{2,:}\\(XH)_{2,:}\end{pmatrix}
 &=-14(22z+19).
\end{align}
They remain nonzero near the critical point.  The disconnected-support
argument is unchanged: the two fixed slow residues determine their pair,
and the fast singleton has exits $(z+\delta,z-\delta)$.

The envelope coefficients~\eqref{eq:fcoeff} were derived for independent
input heights.  Their discriminant now obeys the exact identity
\begin{equation}
 \Delta(z,\delta):=f_1^2-4f_0f_2
 =-\left(\frac{11}{1200}+\frac{11\sqrt6}{4200}\right)
 \left(z+\frac{19}{22}\right)^2\frac{\mathcal P_\delta(z)}{352}.
 \label{eq:perturbeddisc}
\end{equation}
For completeness, the accompanying rational interval certificate verifies
the strict reductions and construction margins on the entire rectangle
\begin{equation}
 \mathcal R=[10.55,10.565]\times[-0.005,0.005].\label{eq:unfoldrect}
\end{equation}
In particular $A>B_0>2$, $a>7/10$,
\begin{equation}
 A-\frac{L}{a}\left[m+\frac{B_0+1}{R_w}\right]>6.691,
 \qquad
 \frac{L(m+1)}{L+R_w}-\frac1{k-R_w/A}>0.0182.
\end{equation}
These are uniform reductions for every candidate triangle: they exclude
negative bottom abscissas and place the bottom to the right of the wedge
apex.  The same envelope maximizations therefore apply to all connected
supports.  The certificate also gives $-f_0>0.102$, $f_1>5.267$,
$-f_2>14.445$, and strictly positive versions of every nonactive margin
listed in Table~\ref{tab:margins}.  Its exact rational endpoints, rather
than these shortened decimal summaries, establish the signs.

Furthermore $\partial_z\mathcal P_\delta>486081$ on $\mathcal R$, and
$\mathcal P_\delta$ has opposite signs at its two $z$ endpoints uniformly
in $\delta$.  Thus each $|\delta|\le0.005$ has one critical root in the
rectangle.  Nonzero $\partial_z\mathcal P_\delta$ gives a smooth root
curve.  At $\delta=0$ this also follows from
\begin{equation}
 \mathcal P'(10+x)=1408x^3+32736x^2+223032x+353380>0
 \quad(x\ge0),
\end{equation}
which proves~\eqref{eq:perturbedslope} without a floating-point rank test.

Above the curve $\Delta<0$, so the envelope is negative everywhere and
all connected supports are excluded.  At the curve its sole zero is
$t_c=-f_1/(2f_2)>0$.  Equality fixes one triangle and its unique positive
normalization.  The active input/output and tangent-face equalities
delete the same three reciprocal pairs as in the balanced endpoint;
the strictly positive remaining margins prove existence and positivity
of all remaining rates.  Below the curve, $\Delta>0$ and
$F(t_c)=-\Delta/(4f_2)>0$.  The three successive vertex changes in
Appendix~\ref{app:subcritical} again give a complete positive realization.
Its exact-law hidden-pair shear makes entropy unbounded.

This proves the corresponding generator statements on $\mathcal R$.
At the critical curve the source and the path have fixed reciprocal
supports and strictly positive present rates, so their stationary laws
and entropy rates depend continuously on $\delta$.  The strictly positive
certified entropy gap at $\delta=0$ therefore remains positive in a
possibly smaller neighborhood.  This establishes the stated two-value
entropy assertion; the geometry certificate alone is not an interval
certificate for that entropy gap on all of $\mathcal R$.
\end{proof}
